% !PRINTED 281
\weberpage{281}

%<<CONT
from which one obtains, by composition,

\weq{16}{\psi' \Theta' \chi' = \begin{pmatrix} 0, & 1 \\ -1, & 0 \end{pmatrix}}

We may accordingly, in the three cases, take the following generating substitutions of the tetrahedral, octahedral and icosahedral group [(6), (13), (15)]:

\[
p = 5. \quad \begin{pmatrix} 2, & 0 \\ 0, & 3 \end{pmatrix}, \ \begin{pmatrix} 0, & 1 \\ -1, & 0 \end{pmatrix}, \ \begin{pmatrix} 2, & 2 \\ 1, & -1 \end{pmatrix}
\]
\[
p = 7. \quad \begin{pmatrix} 3, & 0 \\ 0, & 5 \end{pmatrix}, \ \begin{pmatrix} 0, & 1 \\ -1, & 0 \end{pmatrix}, \ \begin{pmatrix} 1, & -2 \\ -1, & 3 \end{pmatrix}
\]
\[
p = 11. \quad \begin{pmatrix} 3, & 0 \\ 0, & 4 \end{pmatrix}, \ \begin{pmatrix} 0, & 1 \\ -1, & 0 \end{pmatrix}, \ \begin{pmatrix} -3, & 1 \\ 1, & 3 \end{pmatrix}.
\]

By repeated composition, which can be arranged in very many different ways, the complete groups are easily obtained from these:

\[
\begin{gathered}
p = 5.\\[2pt]
\begin{pmatrix} x, & 0 \\ 0, & x^{-1} \end{pmatrix}, \ \begin{pmatrix} 0, & x \\ -x^{-1}, & 0 \end{pmatrix}, \ \begin{pmatrix} x, & y \\ 2y^{-1}, & 3x^{-1} \end{pmatrix} \quad \begin{aligned} x &= 1, 2\\ y &= \pm 1, \pm 2 \end{aligned}
\end{gathered}
\]
\[
\begin{gathered}
p = 7.\\[2pt]
\begin{pmatrix} x, & 0 \\ 0, & x^{-1} \end{pmatrix}, \ \begin{pmatrix} 0, & x \\ -x^{-1}, & 0 \end{pmatrix}, \ \begin{pmatrix} x, & -xz \\ -2x^{-1}z^{-1}, & 3x^{-1} \end{pmatrix}, \ \begin{pmatrix} xz, & x \\ -3x^{-1}, & -2x^{-1}z^{-1} \end{pmatrix}\\[2pt]
x = 1, 2, 3; \quad z = 1, 2, 4\\
(z \text{ quadratic residue of } 7).
\end{gathered}
\]
\[
\begin{gathered}
p = 11.\\[2pt]
\begin{pmatrix} x, & 0 \\ 0, & x^{-1} \end{pmatrix}, \ \begin{pmatrix} 0, & x \\ -x^{-1}, & 0 \end{pmatrix}, \ \begin{pmatrix} x, & -xz \\ -x^{-1}z^{-1}, & 2x^{-1} \end{pmatrix}, \ \begin{pmatrix} xz, & x \\ -2x^{-1}, & -x^{-1}z^{-1} \end{pmatrix}\\[2pt]
x = 1, 2, 3, 4, 5; \quad z = 1, 3, 4, 5, 9\\
(z \text{ quadratic residue of } 11).
\end{gathered}
\]

The application of the transformation (14) leads, for $p = 5$, to no new group; for $p = 7, 11$ one obtains in each case one other group, which arises from these when one puts for $z$ the sequence of the quadratic non-residues instead of the residues\textsuperscript{1}).

\bigskip
\hrule
\bigskip

\textsuperscript{1}) The property of the congruence groups of being simple was already known to \gesperrt{Galois}. Likewise the divisors of index $p$ for $p = 5, 7, 11$ were known to him. These groups have been investigated more thoroughly by
%CONT>>
